% !TEX root = ../main.tex
\section{Conclusiones grupales}
La implementación de la Guía 07 permitió relacionar los fundamentos algebraicos con resultados verificables. La división de prueba ofrece certeza, pero su límite de aproximadamente $2^{L/2}$ divisiones para entradas de $L$ bits impide aplicarla a tamaños criptográficos. Fermat reduce el costo mediante exponenciación modular, aunque su aceptación no demuestra primalidad: $341$ supera la base $2$, y los números de Carmichael $561$, $1105$ y $1729$ superan todas las bases coprimas. Por ello, el equipo separó explícitamente los resultados ``compuesto'', ``primo certificado'' y ``probable primo''.

Miller--Rabin mostró una ventaja esencial: la cadena obtenida de $n-1=2^sd$ puede hallar testigos que Fermat no encuentra. En particular, la base $2$ rechaza a $341$ al descubrir la raíz no trivial $32$ de $1$ módulo $341$; los máximos comunes divisores con $31$ y $33$ recuperan sus factores. Las trazas registran parámetros, residuos y operaciones reales, haciendo auditables las decisiones. Comprobamos las propiedades con una suite de 32 pruebas, una criba independiente para enteros pequeños y 1000 enteros aleatorios menores que un millón contrastados con división de prueba.

Los experimentos recorren siete longitudes entre 16 y 1024 bits y producen 210 probables primos, con candidatos, descartes por precribado, rondas, cuadrados y tiempos registrados. La comparación usa las mismas entradas y bases preseleccionadas para Fermat y Miller--Rabin, evitando atribuir al azar diferencias de rendimiento. Los candidatos de 16 y 32 bits también se certificaron con división de prueba; los mayores mantienen la condición de probables primos. Medir tiempo y contar operaciones son tareas distintas: Python maneja enteros de longitud creciente y su costo no depende solo del número de rondas.

Finalmente, una tasa empírica de error nula no implica imposibilidad de error. La cota $4^{-r}$ exige bases independientes uniformes, y una semilla fija es útil para reproducibilidad, pero inapropiada para claves. Un precribado correcto ahorra exponenciaciones sin transformar por sí solo una prueba probabilística en determinista. El producto entregado conserva código, pruebas, CSV, gráficos e informe reproducibles para permitir repetir y discutir cada conclusión.
